\chapter{Evaluation}
\label{ch:eval}

This chapter tests the central claim of the thesis in five steps. We first
establish how much model \emph{binding} actually needs, with routing held fixed
(Section~\ref{sec:p1}). We then explain \emph{why} a small free-form model fails
(Section~\ref{sec:p2}), and show that the value of the reflexer depends on the
difficulty of the device interface (Section~\ref{sec:p3}). Section~\ref{sec:p4}
reports the result the rest of the thesis turns on: there are task difficulties
that model scale provably cannot solve, but a distilled pattern can.
Section~\ref{sec:p5} closes the loop by distilling those patterns automatically
and measuring how much the choice of teacher matters. The last three sections
place the work in context: real-device coverage (Section~\ref{sec:coverage}),
the honest bottleneck (Section~\ref{sec:routing}), and the economics
(Section~\ref{sec:economics}).

Every number below is taken from a saved run under \texttt{benchmark/reports/};
the cited filename accompanies each result. Two evaluation disciplines hold
throughout, and are detailed in Chapter~\ref{ch:method}: \emph{binding-only}
scoring (routing assumed perfect, so that only the execution stage is under test)
and the \emph{end-state oracle} (success is the final device state, not the
command string). No number in this chapter is estimated or hand-edited.

\section{Choosing the reflexer model}
\label{sec:p0}

Before comparing tiers, the reflexer model itself must be fixed, since the rest of
the chapter holds it constant. The work began with \texttt{ministral-3-3b}, an
instruction-tuned 3B model, which already bound the deterministic slice well; the
question was whether a smaller or different model could do as well or better. We
swept the Qwen3.5 family (0.8B, 2B, 4B) against \texttt{ministral-3-3b} on the
409-atom expansion set, binding-only, each model handed the gold pattern
(\texttt{binding-only-qwen3.5-*-atoms.txt}, \texttt{binding-only-all-ministral-2026-06-16.txt}).

\begin{table}[t]
\centering
\caption{Reflexer model sweep, 409 expansion atoms, binding-only. The 2B is the
smallest model with perfect binding \emph{and} perfect vocabulary compliance.}
\label{tab:reflexer-sweep}
\begin{tabular}{lcc}
\toprule
Model (Q4) & Atom binding & Signature failure mode \\
\midrule
qwen-3.5-0.8b & 380/409 (92.9\%) & violates vocabulary (emits \texttt{docker}, not the \texttt{dock} alias; 0\% docker) \\
\textbf{ministral-3-3b} & 391/409 (95.6\%) & prunes static flags (\texttt{-{}-skip-download}; youtube 81\%) \\
\textbf{qwen-3.5-2b} & \textbf{409/409 (100\%)} & none \\
qwen-3.5-4b & 100\% on captured subset$^{\dagger}$ & none observed \\
\bottomrule
\end{tabular}

\smallskip
{\footnotesize $^{\dagger}$ A full 409-run for 4B was not persisted (only a
docker+service subset, 24/24); it is not reported as a full figure. On the larger
627-task set including core skills, ministral-3-3b scores 601/627 (95.9\%).}
\end{table}

Two findings come out of the sweep. First, \texttt{qwen-3.5-2b} is the best
reflexer found: perfect binding across all 409 atoms, beating both the smaller
0.8B and the larger ministral-3-3b, at the smallest size that fits a 4\,GB device
budget, so it is adopted for the rest of the thesis. Second, and more
interesting, the small models fail in \emph{model-specific signatures} rather than
a single universal weakness. The 0.8B is \emph{semantically} correct but
\emph{schema}-noncompliant: it emits the real binary name (\texttt{docker start})
instead of the skill's defined alias (\texttt{dock start}), failing every docker
task even at temperature $0.2$, because it cannot override a strong lexical prior.
The 2B complies perfectly, at parity with 4B. Vocabulary compliance is therefore a
\textbf{capability threshold between 0.8B and 2B} rather than a smooth gradient,
and scale beyond 2B buys nothing here. The ministral-3-3b signature is different
again: it obeys the vocabulary perfectly but prunes semantically-secondary static
flags (\texttt{-{}-skip-download}), which is why a capable 3B nonetheless drops to
81\% on the affected skill. The lesson for the design is that a reflexer must be
chosen for \emph{schema obedience}, not raw capability, and that a 2B model already
clears that bar.

\section{The binding spine: distillation substitutes for scale}
\label{sec:p1}

The first experiment fixes routing and asks a single question: given that the
right intent has been identified, how large must a model be to emit the right
command? We compare three conditions on the same matched task set (the
\emph{balanced-10}: one to two tasks per skill across all seven skills):
a 2B reflexer handed the distilled pattern, and a free-form thinker (given only
the skill documentation) at 2B and at 35B.

\begin{table}[t]
\centering
\caption{Binding accuracy on the matched balanced-10 set, with routing held
fixed. The 2B reflexer with a pattern equals the 35B free-form thinker; the same
2B without the pattern reaches only 30\%.}
\label{tab:spine}
\begin{tabular}{lcccc}
\toprule
Condition & Binding & Model size & Median latency & Tokens/task \\
\midrule
Reflexer + pattern & \textbf{10/10 (100\%)} & 2B & single call & n/a (1 call) \\
Thinker, no pattern & \textbf{10/10 (100\%)} & 35B & 6.4 s & $\sim$1.3k \\
Thinker, no pattern & 3/10 (30\%) & 2B & 1.7 s & $\sim$0.7k \\
Thinker, no pattern & 1/10 (10\%) & 4B$^{\dagger}$ & 2.8 s & $\sim$0.7k \\
\bottomrule
\end{tabular}

\smallskip
{\footnotesize $^{\dagger}$ The 4B is configured at context 4000 versus the
35B's 131k, so its thinking crowds the skill document; it is not a clean
capability point. Sources: \texttt{binding-only-phase2-qwen2b-balanced10.txt},
\texttt{thinker-binding-phase2-qwen35b-balanced10.txt},
\texttt{thinker-binding-phase2-qwen2b-balanced10.txt},
\texttt{thinker-binding-phase2-qwen4b-balanced10.txt}.}
\end{table}

\begin{figure}[t]
  \centering
  \includegraphics[width=0.72\textwidth]{fig_binding.pdf}
  \caption{Binding accuracy by condition (balanced-10). Scale alone is
  non-monotonic (4B $<$ 2B $<$ 35B); the pattern takes a 2B to the 35B ceiling.}
  \label{fig:binding}
\end{figure}

The result, stated plainly: \textbf{a 2B reflexer with a distilled pattern
matches a 35B free-form thinker (both 100\%), while the same 2B without the
pattern manages 30\%}. The pattern is worth roughly $17\times$ the parameters on
this task. The effect persists at scale on the full set: the 2B reflexer binds
\textbf{72/72 (100\%)} atoms across every skill, where the 2B thinker manages
only \textbf{14/72 (19.4\%)} (\texttt{binding-only-phase2-qwen2b.txt},
\texttt{thinker-binding-phase2-qwen2b.txt}).

Two readings must travel with this headline. First, scale is an \emph{unreliable}
lever for binding: the 4B scores below the 2B, partly a context confound but
partly a real signal, in that the 4B emits higher-quality \emph{real} Linux
commands (\texttt{nmcli radio wifi off}, \texttt{playerctl pause}) because it knows
the true APIs better, and so obeys the skill's defined surface \emph{less}.
Capability to know and obedience to the surface can trade off. Second, this is a
systems and efficiency result, not a claim that ``2B is smarter than 35B.'' The
reflexer is handed the command template and only fills slots, a strictly easier
task than free-form generation. The legitimate claim is precise: \emph{on the
deterministic slice where intents repeat often enough to distill, a 2B constrained
executor matches a 35B free-form thinker's binding accuracy at a fraction of the
size, latency, and tokens.}

\section{Why the small free-form model fails}
\label{sec:p2}

The 30\% of the unaided 2B thinker is not uniform incompetence; it decomposes
into two distinct and opposite failure modes, visible when the per-skill numbers
are examined (\texttt{thinker-binding-phase2-qwen2b.txt}): app 0/6, browser
0/16, call 0/8, clock 1/14, focus 1/6, media 3/10, and device\_settings 9/12.

\paragraph{Under-triggering (no action at all).} On all eight \texttt{call}
tasks the 2B thinker emitted \emph{zero} commands. Short social imperatives such as
``call mom,'' ``hang up,'' and ``answer'' hit the model's conversational prior, and
it replies in prose instead of invoking a tool. With no command emitted, the task
is an automatic miss. The failure is not wrong output; it is the absence of
output.

\paragraph{Wrong vocabulary (real-shell substitution).} On \texttt{app} tasks the
model \emph{does} act, but free-forms from its pretraining prior: \texttt{spotify
\&}, \texttt{killall firefox}, \texttt{ps aux | grep}, never the skill's defined
\texttt{app open -{}-name}. Capability is present; obedience to the defined
surface is absent.

The diagnostic skill is \texttt{clock}, whose documentation deliberately exposes
its real backend (\texttt{systemd-run}, \texttt{systemctl}). Here even the 35B
thinker binds only 30\% (\texttt{thinker-binding-phase2-qwen35b-first10.txt}),
and the 2B only 7\% (1/14), while the reflexer binds 100\% (14/14). A reasoning
model reads the leaked mechanism and prefers the concrete tool it has a strong
prior for (\texttt{systemctl -{}-user list-timers}) over the abstract
\texttt{clock} alias it has a weak prior for. For the thinker, the skill document
is a \emph{soft} prompt mediated by priors and sampling; for the reflexer, the
\texttt{args\_template} is a \emph{hard} constraint with no prose to reason
around. Both failure modes are erased by the pattern, which hard-codes the
command shape. This is also a productization argument for distillation: the
reflexer is robust to documentation quality, whereas the thinker is not.

\section{The spectrum: efficiency on clean APIs, accuracy on realistic ones}
\label{sec:p3}

Whether the reflexer's benefit shows up as \emph{efficiency} or as \emph{accuracy}
depends on the interface. Two settings benchmarks bracket the range. SET-1 is a
clean, friendly settings API (\texttt{settings set connectivity.wifi off}); SET-2
is a realistic Android API (\texttt{settings put <ns> <key> <value>}) with three
namespaces, non-obvious encodings (brightness 0--255, timeouts in milliseconds,
integer enums for Do-Not-Disturb and location, a string-float font scale), and a
permissive \texttt{put} that silently ``succeeds'' on wrong inputs, exactly how
\texttt{adb shell settings} behaves.

On the \emph{clean} API the small model is already a competent executor when it
is given a deep tool (Table~\ref{tab:set1}). Manual state editing collapses to
about a third; the same model driving the deep \texttt{settings} script reaches
100\%, including two- and three-step compound tasks, because the script rejects
wrong guesses and the model recovers. The gain is real, but it is
\emph{efficiency and competence} (fewer calls, fewer tokens, recovery from
error), not new accuracy the model lacked.

\begin{table}[t]
\centering
\caption{SET-1 (clean API), qwen-3.5-2b. The deep tool is an efficiency win:
$+65$ percentage points, at 2 calls instead of 9. Sources:
\texttt{set1-e1e2-qwen2b.txt}, \texttt{set1-e2-hard-qwen2b.txt}.}
\label{tab:set1}
\begin{tabular}{lccc}
\toprule
Mode & Oracle pass & Median calls & Median tokens \\
\midrule
E1, manual JSON edit (no script) & $\sim$30--35\% & 9 & 3098 \\
E2, \texttt{settings} script & \textbf{100\%} (20/20; hard 10/10) & \textbf{2} & \textbf{775} \\
\bottomrule
\end{tabular}
\end{table}

On the \emph{realistic} API the picture inverts. A free-form 2B thinker over the
48 SET-2 tasks scores \textbf{12/48 (25\%)} (\texttt{set2-thinker-qwen2b-48.txt}).
Its failures cluster exactly on what cannot be guessed: every brightness percent
that must become 0--255, every timeout that must become milliseconds, the integer
enums, and the key-name traps (\texttt{airplane mode} $\rightarrow$
\texttt{airplane\_mode\_on}). It passes only guessable booleans and trivial
arithmetic. Here the abstraction is no longer about efficiency: a small free-form
agent simply \emph{cannot} drive a realistic device-settings API. The reflexer's
benefit has shifted from saving calls to supplying accuracy.

\begin{figure}[t]
  \centering
  \includegraphics[width=0.95\textwidth]{fig_spectrum.pdf}
  \caption{The reflexer's benefit scales with discovery difficulty: an efficiency
  win on a clean API, an accuracy win on a realistic one.}
  \label{fig:spectrum}
\end{figure}

\section{The limit of scale}
\label{sec:p4}

If a small free-form model fails SET-2, the obvious response is to use a bigger
one. This section shows where that response stops working, the result the thesis
is built around.

A 35B reasoning thinker (\texttt{qwen3.6-35b}) was run on the hardest SET-2
slices: the last three tasks (\textbf{3/3}, \texttt{set2-thinker-qwen3.6-35b-*.json})
and the five hardest excluding those (\textbf{2/5}, \texttt{set2-thinker-35b-hard5.txt}),
for \textbf{5/8 (62.5\%)} combined, at roughly nine calls and 17--20 seconds per
task, off-device. The 35B clearly beats the 2B, but the manner of its failures
is the point. They are not \emph{discovery} failures. The 35B reliably finds the
right namespace and key by exploring the API (\texttt{list}, \texttt{grep},
\texttt{get}, then \texttt{put}). Its failures are \emph{encoding semantics}
(Table~\ref{tab:limit}).

\begin{table}[t]
\centering
\caption{The 35B thinker's SET-2 failures are conventions, not discovery. These
values are not derivable from the API; the model guesses, and is wrong.}
\label{tab:limit}
\begin{tabular}{lcc}
\toprule
Task & Correct value & 35B produced \\
\midrule
``largest font'' & \texttt{font\_scale "1.30"} & \texttt{"1.5"} \\
``DND total silence'' & \texttt{zen\_mode 2} & \texttt{1} \\
``DND important only'' & \texttt{zen\_mode 1} & \texttt{2} \\
\bottomrule
\end{tabular}
\end{table}

This separates two kinds of difficulty by what scale can fix.
\textbf{Discovery}, knowing \emph{which} namespace and key, is solvable by
exploration, and the 35B does it. \textbf{Encoding semantics}, knowing what the
enum value \texttt{2} means or what number is ``largest'', is \emph{not
discoverable} from the API at all, and even a 35B guesses and is wrong about half
the time. The \texttt{zen\_mode} error is especially telling: the model makes the
\emph{same} mistake every time it is asked, which marks a consistent gap in what
the model knows rather than sampling noise. \textbf{Scale rescues discovery; scale
cannot rescue undiscoverable conventions.} That knowledge exists only in a
distilled pattern, which is what the next section supplies.

\begin{figure}[t]
  \centering
  \includegraphics[width=0.72\textwidth]{fig_discovery.pdf}
  \caption{Why a distilled pattern beats a $17\times$ larger model on the
  realistic API: the bottleneck is undiscoverable convention, not reasoning.}
  \label{fig:discovery}
\end{figure}

\section{Closing the loop: cold-start distillation and teacher quality}
\label{sec:p5}

The convention knowledge that no model discovers can be supplied once, offline,
by a teacher. We test this in the hardest honest setting: \emph{cold-start}
distillation, in which the teacher is given only the skill documentation and the
API schema (never any task or its gold outcome), and must mint the patterns from
documentation and prior knowledge alone. The 2B reflexer then applies those
patterns on the single-setting tasks, scored by the same end-state oracle.

With a capable cloud teacher (Claude Sonnet 4.6), the reflexer rises from the
free-form 25\% to \textbf{19/25 (76\%)} on the initial cold-start run, and to
\textbf{20/25 (80\%)} under the key-matched mapping used for the teacher
comparison below.% TODO reconcile 19 vs 20: confirm whether one task flipped between the E3/E4 run and the key-matched re-run; report the single authoritative figure.
\ The teacher encoded the very conventions the larger in-context thinker got
wrong (including \texttt{zen\_mode = 2} for total silence), and honestly flagged
the two it was unsure of (\texttt{font\_scale} ``largest'' and
\texttt{wifi\_sleep\_policy} ``never''), which are genuinely device- or
version-specific and not derivable from documentation
(\texttt{set2-reflexer-qwen-3.5-2b-*.json}). A distilled pattern minted without
ever seeing a task therefore beats a $17\times$ larger model that reasons over
the live API, because it carries knowledge the API does not expose.

Because the teacher is a swappable component, its quality should propagate to the
student, and it does. Repeating the cold-start distillation with a \emph{local}
35B teacher (\texttt{qwen3.6-35b}) under an identical prompt, then running the
\emph{same} 2B reflexer on each teacher's patterns, gives Table~\ref{tab:teacher}.

\begin{table}[t]
\centering
\caption{Teacher quality propagates to the student: $-20$ points from the weaker
teacher, on the same reflexer and tasks (key-matched mapping, 25 single-setting
tasks). Sources: \texttt{patterns\_distilled.json} (cloud) vs
\texttt{patterns\_distilled\_pi.json} (local); \texttt{set2-reflexer-qwen-3.5-2b-*}.}
\label{tab:teacher}
\begin{tabular}{lccc}
\toprule
Teacher (distiller) & Structure (keys/ns) & Conventions correct & Reflexer pass \\
\midrule
Sonnet 4.6 (cloud) & 32/32, 0 errors & $\sim$5/7 critical & \textbf{20/25 (80\%)} \\
qwen3.6-35b (local) & 32/32, 0 errors & $\sim$2/7 critical & 15/25 (60\%) \\
\bottomrule
\end{tabular}
\end{table}

Both teachers are perfect on \emph{structure} (every key, every namespace), so
discovery is easy for either. They diverge on \emph{conventions}. The local 35B
baked in its own wrong priors: \texttt{zen\_mode} total-silence as \texttt{1}
(the same error it made as a thinker, a consistent gap), \texttt{location\_mode}
high-accuracy collapsed to a three-value enum, a wrong brightness worked example,
and an over-verbose \texttt{airplane} encoding that induced \emph{collateral
edits} in the student (the 2B then toggled Wi-Fi and Bluetooth too, failing the
no-collateral oracle). It also reported ``all verified'' while silently wrong,
where the cloud teacher flagged its uncertainty. The lesson is operational:
cold-start distillation should use the strongest available teacher as a one-time
step; a weak local teacher is a weak distiller, and the cheap reflexer faithfully
executes whatever priors it inherits. The two conventions both teachers miss
point at the remaining lever, session-based grounding on the real device, which
Chapter~\ref{ch:discussion} returns to.

\section{Coverage: what fraction of real use-cases this addresses}
\label{sec:coverage}

Binding accuracy is meaningful only against a denominator of real demand. Mapping
the skill library against the on-device use-cases that Google (Gemini Nano via
AICore) and Apple (Apple Intelligence, App Intents) target shows that the
reflexer's natural territory is the \emph{deterministic head}: single-command
device and media control, smart-home toggles, system and timer operations, and
parameterized reads with clear verbs. Apple's App Intents are, in effect, the
same construct as a distilled pattern, a pre-declared parameterized action the
on-device model fills rather than invents, which is independent validation of
the executioner design. What the reflexer does not cover, by design, escalates to
the thinker: multi-intent requests, conversation and planning, and anything
outside an installed skill. The qualitative coverage mapping is documented in
\texttt{docs/research/ondevice\_usecases\_pixel\_apple.md}.

\section{The honest bottleneck: routing, not binding}
\label{sec:routing}

This thesis has held routing fixed in order to isolate binding, and binding is
effectively solved on the deterministic slice (96--100\%). End-to-end coverage,
however, is the product of routing and binding, and routing is the limiter. On
creative, well-designed skills, routing reaches about 85\%; on overlapping skills
it falls to 66--68\%, so realistic single-intent coverage today is roughly
65--85\%, set almost entirely by the router
(\texttt{routing\_ablation\_report.md}; \texttt{findings\_retrieval\_and\_selection.md}).
Decomposed, retrieval is not the bound (the correct candidate is usually in the
top-3); the bound is \emph{selection}: the router picks the wrong sibling within
or across skills, and selection does not improve with selector scale, only with
better prompting. Stating this plainly is part of the contribution: the binding
ceiling is reached, and the open problem for a shippable system is the router.

\section{Economics: when reflexing is worth it}
\label{sec:economics}

The architecture is ultimately a cost argument, so it should be settled as cost
accounting. A pattern is a frozen answer shape: the reflexer never \emph{solves}
the task, it fills in a solution found once, which is why a 2B can match a 35B on
it. In distillation terms the large model is the \emph{teacher} that mints the
pattern once, offline, and the reflexer is the \emph{student} that runs it
forever, on-device; the teacher's cost is paid at distillation time, not at
inference time, and, unlike a fine-tune, the symbolic pattern stays
plug-and-play.

\begin{table}[t]
\centering
\caption{Order-of-magnitude cost model. Break-even is reached at roughly ten
invocations, before counting the size, latency, and on-device-feasibility
advantages. Source: \texttt{findings\_thinker\_vs\_reflexer.md} \S9.}
\label{tab:economics}
\begin{tabular}{lll}
\toprule
 & One-time (distill) & Per-invocation (forever) \\
\midrule
Reflexer atom & $\sim$5--15k teacher tokens + brief review & $\sim$0.3k tokens, $<$1 s, fits on-device (2B) \\
Thinker (big) & 0 & $\sim$1.3k tokens, $\sim$6.4 s, off-device (35B) \\
\bottomrule
\end{tabular}
\end{table}

With per-invocation savings of roughly 1k tokens and a one-time cost near 10k
tokens to mint an atom, distillation breaks even at about ten invocations
(Table~\ref{tab:economics}), before the $\sim$17$\times$ per-token size discount,
the drop from 6.4\,s to under a second, and the fact that the reflexer fits
on-device while the 35B does not. For a head intent invoked thousands of times
this is a $100$--$1000\times$ net win; for a tail intent invoked once it never
amortizes and should be routed to the thinker. The right design is therefore not
``reflexer or bigger model or better router'' but all three in their places:
distill the high-frequency head, spend the big model as the offline teacher and
the tail fallback, and fund the router regardless, because it gates both tiers.

\section{Summary}
\label{sec:eval-summary}

The five results compose into one line. A small model cannot drive a realistic
device API by reasoning (25\%), and scaling the model up to 35B rescues
\emph{discovery} but not \emph{convention} (62.5\%, with consistent
encoding errors). A pattern distilled once by a strong teacher carries exactly
that convention, lifting a 2B reflexer to 76--80\% at one on-device call, and the
quality of that teacher propagates directly to the student ($80\%$ vs $60\%$). On
the clean API the same abstraction is instead an efficiency win
($+65$ points, 2 calls versus 9). Binding, in short, is solved by distillation
rather than scale; the remaining open limit is routing. The next chapter draws
out what this means for design, for productization, and for the validity of the
claims.
